\documentclass[11pt]{article}

\usepackage[margin=1.02in]{geometry}
\usepackage{amsmath,amssymb,amsthm,mathtools}
\usepackage{booktabs,array}
\usepackage{microtype}
\usepackage[hidelinks]{hyperref}

\allowdisplaybreaks
\setlength{\parindent}{1.2em}
\setlength{\parskip}{0.15em}

\newtheorem{theorem}{Theorem}[section]
\newtheorem{lemma}[theorem]{Lemma}
\newtheorem{proposition}[theorem]{Proposition}
\newtheorem{corollary}[theorem]{Corollary}
\theoremstyle{definition}
\newtheorem{remark}[theorem]{Remark}

\newcommand{\per}{\operatorname{per}}
\newcommand{\pcap}{\operatorname{Cap}}
\newcommand{\e}{\mathrm e}
\newcommand{\Rge}{\mathbb R_{\ge0}}

\title{A Penalized Core--Entropy Proof of Dittert's Conjecture in Dimension Six}
\author{Extension note to the dimensions $7$--$15$ argument}
\date{24 July 2026}

\begin{document}
\maketitle

\begin{abstract}
This note gives a proposed completion of the Hall-deficit proof of Dittert's
conjecture in dimension six.  Two losses that are harmless in dimension seven
become decisive in dimension six.  First, the sharp zero-rectangle permanent
bound is stabilized by the nonuniformity of the critical-core marginals.  For
general codimension, differentiating one zero-block column with its exact
entropy factor gives
\[
 \per(B)\ge L^{\sharp}_{n,a,b}
 \exp\!\left\{\frac{s_I+s_J}{2(a^2+b^2)}\right\}.
\]
In codimension one, the exact cofactor entropy from the dimension-seven note
has stronger dimension-specific quadratic moduli.  Second, in the critical
core stationarity equation we retain both terms discarded by the earlier
scalar reduction: the linear conditional-product loss and the exponential
core-entropy loss.  Eliminating the two auxiliary energy variables yields the
penalized algebraic criterion
\[
 4\lambda\theta(\mathcal B+hR)(A_0\mathcal B+qR)>R^2.
\]
All two-sided Hall configurations, including both mixed-sign quadrants, are
then covered by exact bivariate Bernstein subdivision.  The one-sided cuts are
excluded by the existing entropy dichotomy with an improved exact row-energy
coefficient.  The accompanying standard-library Python script uses rational
arithmetic only.
\end{abstract}

\section{Statement and inherited facts}

For a nonnegative $n\times n$ matrix $A=(a_{ij})$ whose entries sum to $n$,
write $r_i$ and $c_j$ for its row and column sums and put
\[
 \Phi(A)=\prod_i r_i+\prod_jc_j-\per(A),
 \qquad
 \gamma_n=\frac{n!}{n^n}.
\]
We use the arbitrary-dimensional parts of the accompanying manuscripts:

\begin{enumerate}
\item every entrywise positive global maximizer is $J_n/n$;
\item a hypothetical boundary maximizer has deficits
\[
 P:=\prod_i r_i=1-\rho,
 \qquad Q:=\prod_jc_j=1-\sigma,
 \qquad H:=\per(A)=\gamma_n-\delta
\]
with
\[
 0\le\rho,\sigma,
 \qquad \rho+\sigma\le\delta\le\gamma_n;
 \tag{1.1}\label{eq:budget}
\]
\item a critical Hall pair gives a dilation $S=A/(1-\tau)$ dominating a
doubly stochastic matrix $B$ with a forced zero rectangle;
\item if $B$ has an $a\times b$ zero rectangle and
$p=n-a-b\ge1$, then
\[
 \per(B)\ge L^{\sharp}_{n,a,b}
 :=\frac{v_{n-a}v_{n-b}}{v_p},
 \qquad v_k=\frac{k!}{k^k},\quad v_0=1.
 \tag{1.2}\label{eq:sharp}
\]
\end{enumerate}

The result proved below is the missing dimension.

\begin{theorem}\label{thm:n6}
For every nonnegative $6\times6$ matrix $A$ whose entries sum to $6$,
\[
 \Phi(A)\le 2-\frac{6!}{6^6}=2-\frac5{324}.
\]
Equality holds only at $A=J_6/6$.
\end{theorem}

It remains only to exclude a boundary global maximizer.

\section{Entropy stability of a sharp zero rectangle}

\subsection{A one-column refinement in arbitrary codimension}

Let a doubly stochastic $B$ have a zero block $B[R,C]=0$, where
$|R|=a$, $|C|=b$, and $p=n-a-b\ge1$.  Put
\[
 I=R^c,
 \quad J=C^c,
 \quad m=|I|=n-a,
 \quad \ell=|J|=n-b.
\]
The core $X=B[I,J]$ has total mass $p$.  Let
\[
 \xi_i=\sum_{j\in J}x_{ij},
 \qquad
 \eta_j=\sum_{i\in I}x_{ij},
\]
and define
\[
 s_I=\sum_{i\in I}\left(\xi_i-\frac pm\right)^2,
 \qquad
 s_J=\sum_{j\in J}\left(\eta_j-\frac p\ell\right)^2.
 \tag{2.1}\label{eq:sIJ}
\]

\begin{theorem}[One-column stability of the sharp block bound]
\label{thm:one-column-stability}
Under the preceding assumptions,
\[
 \per(B)\ge L^{\sharp}_{n,a,b}
 \exp\!\left\{
 \frac12\max\left(\frac{s_I}{b^2},\frac{s_J}{a^2}\right)
 \right\}.
 \tag{2.2}\label{eq:max-entropy}
\]
Consequently,
\[
 \per(B)\ge L^{\sharp}_{n,a,b}
 \exp\!\left\{\frac{s_I+s_J}{2(a^2+b^2)}\right\}.
 \tag{2.3}\label{eq:general-block-entropy}
\]
\end{theorem}

\begin{proof}
Choose a column $c\in C$ and write its nonzero part as the probability vector
$y=(y_i)_{i\in I}$.  Let
\[
 P_B(x)=\prod_{r=1}^n\left(\sum_{j=1}^n b_{rj}x_j\right),
 \qquad
 q=\left.\frac{\partial P_B}{\partial x_c}\right|_{x_c=0}.
\]
The exact first step in the column-entropy argument gives
\[
 \pcap(q)\ge F(y),
 \qquad
 F(y):=\prod_{i\in I}(1-y_i)^{1-y_i}.
 \tag{2.4}\label{eq:first-column-cap}
\]
Differentiate next in the remaining $b-1$ variables indexed by $C$.  Their
successive degrees are at most
$m-1,m-2,\ldots,p+1$.  The derivative-capacity inequality supplies
\[
 \frac{v_{m-1}}{v_p}.
\]
After these differentiations, square-free coefficient extraction in the
remaining $\ell$ variables supplies $v_\ell$.  Therefore
\[
 \per(B)\ge F(y)\frac{v_{m-1}v_\ell}{v_p}
 =L^{\sharp}_{n,a,b}\frac{F(y)}{G(m)},
 \tag{2.5}\label{eq:selected-column-bound}
\]
where $G(m)=((m-1)/m)^{m-1}$.

For $h(t)=(1-t)\log(1-t)$, one has $h''(t)=1/(1-t)\ge1$.  Since
$\sum_i y_i=1$,
\[
 \log\frac{F(y)}{G(m)}
 =\sum_{i\in I}h(y_i)-m h(1/m)
 \ge\frac12\sum_{i\in I}\left(y_i-\frac1m\right)^2.
 \tag{2.6}\label{eq:column-strong-convexity}
\]

Average the $b$ columns in $C$.  Their average on $I$ is
\[
 \bar y_i=\frac{1-\xi_i}{b},
 \qquad
 \bar y_i-\frac1m=-\frac1b\left(\xi_i-\frac pm\right).
\]
By convexity of squared Euclidean distance, at least one column satisfies
\[
 \sum_i\left(y_i-\frac1m\right)^2\ge\frac{s_I}{b^2}.
\]
Equations \eqref{eq:selected-column-bound} and
\eqref{eq:column-strong-convexity} give the first term inside the maximum in
\eqref{eq:max-entropy}.  Apply the same argument to $B^{\mathsf T}$ for the
second term.  Finally,
\[
 \max\left(\frac{x}{b^2},\frac{y}{a^2}\right)
 \ge\frac{x+y}{a^2+b^2}
 \qquad(x,y\ge0),
\]
which proves \eqref{eq:general-block-entropy}.
\end{proof}

\subsection{Stronger codimension-one moduli}

When $p=1$, the cofactor argument from the dimension-seven note gives more
than \eqref{eq:general-block-entropy}.  If $m=b+1$ and $\ell=a+1$, it gives the
exact inequality
\[
 \per(B)\ge L^{\sharp}_{n,a,b}
 \exp\{\mathcal E_m(\xi)+\mathcal E_\ell(\eta)\},
 \tag{2.7}\label{eq:exact-p1-entropy}
\]
where, for a probability vector $z\in\mathbb R^k$,
\[
 \mathcal E_k(z)=
 \sum_{i=1}^k(1-z_i)\log(1-z_i)
 -k\left(1-\frac1k\right)\log\left(1-\frac1k\right).
 \tag{2.8}\label{eq:Ek}
\]
The following small-dimensional constants are the useful improvement over the
uniform coefficient $1/2$.

\begin{lemma}[Pointwise entropy moduli]\label{lem:small-entropy-moduli}
For every probability vector $z$,
\[
 \mathcal E_k(z)\ge c_k
 \sum_i\left(z_i-\frac1k\right)^2,
 \tag{2.9}\label{eq:entropy-modulus}
\]
where one may take
\[
 c_2=c_5=\frac{23}{40},
 \qquad
 c_3=c_4=\frac{301}{500}.
 \tag{2.10}\label{eq:ck}
\]
\end{lemma}

\begin{proof}
Put $h(t)=(1-t)\log(1-t)$, $u=1/k$, and
\[
 g_{k,c}(t)=h(t)-h(u)-h'(u)(t-u)-c(t-u)^2.
\]
It is enough to prove $g_{k,c}(t)\ge0$ on $[0,1]$, because the linear terms
cancel after summation.  Now
\[
 g_{k,c}''(t)=\frac1{1-t}-2c.
\]
For every pair $(k,c)$ in \eqref{eq:ck}, the unique sign change of
$g''$ occurs strictly before $u$.  Thus $g$ is concave first and then convex;
using $g(u)=g'(u)=0$, its only possible additional minimum is the endpoint
$t=0$.  At that endpoint,
\[
 g_{k,c}(0)=\log\frac{k}{k-1}-\frac1k-\frac{c}{k^2}.
 \tag{2.11}\label{eq:g0}
\]
The required four inequalities are
\[
 \log\frac54>\frac{223}{1000},
 \quad
 \log2>\frac{103}{160},
 \quad
 \log\frac43>\frac{2301}{8000},
 \quad
 \log\frac32>\frac{1801}{4500}.
\]
They follow, respectively, from the even alternating partial sums
\[
 \frac{27419}{122880},
 \qquad \frac23,
 \qquad \frac{12581}{43740},
 \qquad \frac{77}{192}.
\]
This proves the pointwise inequality and hence \eqref{eq:entropy-modulus}.
\end{proof}

For the two codimension-one pairs in dimension six, Lemma
\ref{lem:small-entropy-moduli} and \eqref{eq:exact-p1-entropy} give
\[
 \per(B)\ge L^{\sharp}_{6,1,4}\e^{(23/40)(s_I+s_J)},
 \tag{2.12}\label{eq:lambda14}
\]
and
\[
 \per(B)\ge L^{\sharp}_{6,2,3}\e^{(301/500)(s_I+s_J)}.
 \tag{2.13}\label{eq:lambda23}
\]

\section{Dimension-six localization and inverse energy}

For the rest of the proof,
\[
 n=6,
 \qquad
 \gamma=\gamma_6=\frac5{324},
 \qquad
 \eta=\frac1{27}.
\]
Define $\alpha_0=0$ and
\[
 \alpha_k=
 \begin{cases}
 12\left(\dfrac{k}{6}+\eta\right)
       \left(1-\dfrac{k}{6}-\eta\right),&2k<6,\\[2mm]
 \dfrac{2k(6-k)}6,&2k\ge6.
 \end{cases}
 \tag{3.1}\label{eq:alpha6}
\]
The exact values are
\[
 (\alpha_0,\ldots,\alpha_5)
 =\left(0,\frac{473}{243},\frac{680}{243},3,\frac83,\frac53\right).
 \tag{3.2}\label{eq:alpha-values}
\]
The relative-entropy subset proof applies because
\[
 \eta^2-\frac{\gamma}{12(1-\gamma)}
 =\frac{61}{930204}>0,
 \qquad
 \frac13+\eta<\frac12.
 \tag{3.3}\label{eq:localization6}
\]
In particular, $\alpha_k\le3$.

Let $E$ be any row subset and write its sum as $|E|+e$.  Applying the subset
bound either to $E$ or to its complement gives
\[
 |e|^2\le\frac{3\rho}{1-\delta}
 \le\frac{3\gamma}{1-\gamma}=\frac{15}{319}.
\]
The same is true for column subsets.  Hence every signed Hall excess below
satisfies
\[
 |e|,|f|<T_0,
 \qquad
 T_0=\frac{217}{1000},
 \tag{3.4}\label{eq:signed-box}
\]
because
\[
 T_0^2-\frac{15}{319}=\frac{21391}{319000000}>0.
\]

The dimension-six argument needs a stronger conditional inverse-energy
coefficient than the coarse global estimate.  The following finite ladder
will be used only for the pair $(a,b)=(2,3)$.

\begin{center}
\begin{tabular}{c|c|c}
\toprule
$A_*$&$r_*$&$\theta_*$\\
\midrule
$0$&$169/200$&$7/20$\\
$3/800$&$86533/100000$&$37/100$\\
$7/1250$&$8761/10000$&$19/50$\\
$73/10000$&$8869/10000$&$39/100$\\
$161/20000$&$17843/20000$&$79/200$\\
$17/2000$&$4477/5000$&$199/500$\\
$4383/500000$&$35897/40000$&$2/5$\\
\bottomrule
\end{tabular}
\end{center}

\begin{lemma}[Energy ladder]\label{lem:energy-ladder}
Suppose a permanent lower bound gives $H\ge A_*$, where $(A_*,r_*,\theta_*)$
is one row of the table.  Then every row and column sum is larger than $r_*$.
Moreover, for every row group $I$ with groupwise AM--GM envelope $P_a(e)$,
\[
 V_I:=\sum_{i\in I}\left(
 \frac1{r_i}-\frac1{|I|}\sum_{h\in I}\frac1{r_h}
 \right)^2
 \le\frac{P_a(e)-P}{\theta_*}.
 \tag{3.5}\label{eq:conditional-energy-tier}
\]
The column analogue is identical.
\end{lemma}

\begin{proof}
The budget and $H=\gamma-\delta\ge A_*$ give
\[
 P,Q\ge1-\delta\ge1-\gamma+A_*=:P_*.
 \tag{3.6}\label{eq:Pstar}
\]
If one row sum were at most $r_*<1$, AM--GM on the other five rows would give
\[
 P\le r_*\left(\frac{6-r_*}{5}\right)^5<P_*.
\]
For every table row, this last strict rational inequality and
\[
 \frac{P_*r_*^2}{2}>\theta_*
 \tag{3.7}\label{eq:energy-table-check}
\]
are verified exactly in the accompanying script.

Let $q$ be the mean row sum on $I$.  Since all row sums and $q$ are at least
$r_*$,
\[
 \frac rq-1-\log\frac rq
 \ge\frac{r_*^2}{2}\left(\frac1r-\frac1q\right)^2.
 \tag{3.8}\label{eq:scalar-energy}
\]
The groupwise linear terms cancel, so
\[
 \log\frac{P_a(e)}P\ge\frac{r_*^2}{2}V_I.
\]
Finally,
\[
 P_a(e)-P\ge P\log\frac{P_a(e)}P
 \ge\frac{P_*r_*^2}{2}V_I>\theta_*V_I.
\]
\end{proof}

The first row of the table is unconditional.  In particular, throughout the
proof one may use
\[
 V_I\le\frac{20}{7}\bigl(P_a(e)-P\bigr).
 \tag{3.9}\label{eq:global-conditional-energy}
\]

\section{The penalized critical-core criterion}

Let $I,J$ be a two-sided critical Hall pair.  Put
\[
 R=I^c,
 \quad C=J^c,
 \quad a=|R|,
 \quad b=|C|,
 \quad p=6-a-b\ge1,
\]
and, after transposition, assume $1\le a\le b$.  Define the signed excesses
\[
 e=\sum_{i\in R}r_i-a,
 \qquad
 f=\sum_{j\in C}c_j-b,
 \qquad
 x=s_A(R,C).
 \tag{4.1}\label{eq:efx}
\]
Then
\[
 p\tau=e+f-x.
 \tag{4.2}\label{eq:tau-signed}
\]
Set
\[
 e_+=\max(e,0),
 \qquad f_+=\max(f,0),
 \qquad T=e_++f_+,
 \qquad \omega=1-\frac Tp.
 \tag{4.3}\label{eq:Tomega}
\]
Thus $p\tau\le T$ and $1-\tau\ge\omega$.

\subsection{Core saturation and general stationarity}

The doubly stochastic cut identity gives $s_B(I,J)=p$, while criticality gives
$s_S(I,J)=p$.  Since $B\le S$, the core saturates entrywise:
\[
 B[I,J]=S[I,J]=\frac{A[I,J]}{1-\tau}.
 \tag{4.4}\label{eq:core-saturation-general}
\]
Let $\xi,\eta$ be the core marginals and let $s=s_I+s_J$ as in
\eqref{eq:sIJ}.

For a permutation, the number of selected entries in $I\times J$ exceeds the
number selected in $R\times C$ by exactly $p$.  Let $H_k$ be the permanent
weight of permutations using exactly $k$ entries of $R\times C$; then they use
$p+k$ core entries and $0\le k\le d:=\min(a,b)$.

\begin{lemma}[General critical-core stationarity]\label{lem:general-core-stationarity}
One has
\[
 P\left(\sum_{i\in I}\frac{\xi_i}{r_i}-p\right)
 +Q\left(\sum_{j\in J}\frac{\eta_j}{c_j}-p\right)
 =\frac{p\tau H+\sum_{k\ge1}kH_k}{1-\tau}.
 \tag{4.5}\label{eq:general-stationarity}
\]
\end{lemma}

\begin{proof}
Scale $A[I,J]$ by a parameter and globally renormalize.  The core has total
mass $p(1-\tau)$.  Differentiation of the row product, column product, and
permanent gives, before simplification,
\[
 (1-\tau)P\sum_{i\in I}\frac{\xi_i}{r_i}
 +(1-\tau)Q\sum_{j\in J}\frac{\eta_j}{c_j}
 -\sum_k(p+k)H_k-p(1-\tau)\Phi(A)=0.
\]
Use $\Phi(A)=P+Q-H$ and divide by $1-\tau$.
\end{proof}

Let $H_0$ be the contribution of permutations avoiding $R\times C$.  Every
matching of $B$ is of this type, and $A\ge(1-\tau)B$ entrywise on its support.
Sections 2.1 and 2.2 therefore give
\[
 H_0\ge A_0\e^{\lambda s},
 \qquad
 A_0=L\omega^6,
 \qquad
 L=L^{\sharp}_{6,a,b},
 \tag{4.6}\label{eq:H0-general}
\]
where
\[
 \lambda=
 \begin{cases}
 \dfrac{23}{40},&(a,b)=(1,4),\\[1mm]
 \dfrac{301}{500},&(a,b)=(2,3),\\[1mm]
 \dfrac1{2(a^2+b^2)},&p\ge2.
 \end{cases}
 \tag{4.7}\label{eq:lambda-pairs}
\]
Since $\sum kH_k\le d(H-H_0)$, the right side of
\eqref{eq:general-stationarity} is at most
\[
 \frac{(T+d)H-dH_0}{\omega}.
 \tag{4.8}\label{eq:stationarity-upper-general}
\]

\subsection{Signed product envelopes}

Put
\[
 P_a(e)=\left(1+\frac ea\right)^a
        \left(1-\frac e{6-a}\right)^{6-a},
 \qquad
 P_b(f)=\left(1+\frac fb\right)^b
        \left(1-\frac f{6-b}\right)^{6-b}.
 \tag{4.9}\label{eq:signed-envelopes}
\]
These formulas are valid for either sign in the box \eqref{eq:signed-box}, and
AM--GM gives $P\le P_a(e)$ and $Q\le P_b(f)$.  Define
\[
 u_r=P_a(e)-P,
 \qquad
 u_c=P_b(f)-Q,
\]
\[
 D=2-P_a(e)-P_b(f),
 \qquad
 K=\gamma-D,
 \qquad
 u=\delta-D.
 \tag{4.10}\label{eq:DKu}
\]
The joint budget implies
\[
 u\ge u_r+u_c\ge0,
 \qquad
 H=K-u.
 \tag{4.11}\label{eq:u-general}
\]
Write
\[
 m=6-a,
 \qquad \ell=6-b,
 \qquad
 E=\frac e{m-e},
 \qquad F=\frac f{\ell-f}.
 \tag{4.12}\label{eq:EF-general}
\]
We use $E_+=\max(E,0)$ and $F_+=\max(F,0)$ below.
The harmonic-mean inequality and Lemma \ref{lem:energy-ladder} give
\[
 \sum_{i\in I}\frac{\xi_i}{r_i}-p
 \ge pE-\sqrt{s_I V_I},
 \tag{4.13}\label{eq:harmonic-core-row}
\]
and the analogous column inequality.  If a coefficient $E$ or $F$ is
negative, its product-loss term has the favorable sign.  Thus, with any valid
energy coefficient $\theta$,
\[
 \begin{aligned}
 &P\left(\sum_{i\in I}\frac{\xi_i}{r_i}-p\right)
 +Q\left(\sum_{j\in J}\frac{\eta_j}{c_j}-p\right)\\
 &\qquad\ge
 pP_a(e)E+pP_b(f)F
 -pE_+u_r-pF_+u_c
 -\sqrt{\frac{s(u_r+u_c)}{\theta}}.
 \end{aligned}
 \tag{4.14}\label{eq:stationarity-lower-penalized}
\]

Set
\[
 \mathcal B=
 pP_a(e)E+pP_b(f)F
 -\frac{(T+d)K-dA_0}{\omega},
 \tag{4.15}\label{eq:B-general}
\]
\[
 C_0=\frac{T+d}{\omega},
 \qquad
 h=C_0-p(E_++F_+),
 \qquad
 q=\frac{dA_0}{\omega},
 \qquad
 R_0=K-A_0.
 \tag{4.16}\label{eq:hqR}
\]
On the full box \eqref{eq:signed-box}, $h>0$ for all six unordered Hall pairs.
A simple uniform check is
\[
 h\ge d-\frac{2pT_0}{\ell-T_0}>0.
 \tag{4.17}\label{eq:h-positive}
\]

Put
\[
 v=u_r+u_c,
 \qquad
 w=\e^{\lambda s}-1.
\]
Combining \eqref{eq:stationarity-upper-general} and
\eqref{eq:stationarity-lower-penalized}, while retaining rather than discarding
the two coercive terms, gives
\[
 \mathcal B+qw+hv\le\sqrt{\frac{sv}{\theta}}.
 \tag{4.18}\label{eq:penalized-energy}
\]
The permanent budget gives
\[
 v+A_0w\le u+A_0w\le R_0.
 \tag{4.19}\label{eq:triangle-budget}
\]
Since $s=\lambda^{-1}\log(1+w)\le w/\lambda$, a necessary condition is
\[
 \lambda\theta(\mathcal B+qw+hv)^2\le wv.
 \tag{4.20}\label{eq:quadratic-triangle}
\]

\subsection{Eliminating the two energy variables}

\begin{lemma}[Penalized energy triangle]\label{lem:penalized-triangle}
Let
$A_0,R_0,\mathcal B,q,h,\lambda,\theta$ be nonnegative, with
$A_0,\mathcal B,h,\lambda,\theta>0$.  There do not exist $w,v\ge0$ satisfying
\eqref{eq:triangle-budget} and \eqref{eq:quadratic-triangle} whenever
\[
 4\lambda\theta
 (\mathcal B+hR_0)(A_0\mathcal B+qR_0)>R_0^2.
 \tag{4.21}\label{eq:enhanced-master}
\]
\end{lemma}

\begin{proof}
Put $c=\lambda\theta$ and
\[
 F(w,v)=wv-c(\mathcal B+qw+hv)^2.
\]
The axes have $F<0$.  At a positive interior stationary point, with
$Z=\mathcal B+qw+hv$, one has
\[
 v=2cqZ,
 \qquad w=2chZ,
 \qquad Z(1-4cqh)=\mathcal B>0.
\]
Consequently,
\[
 F=cZ^2(4cqh-1)<0.
\]
Thus a nonnegative maximum over the triangle
$A_0w+v\le R_0$ must lie on its sloping edge.  Substituting
$v=R_0-A_0w$ gives the concave quadratic
\[
 \begin{aligned}
 F(w,R_0-A_0w)
 &=-c(\mathcal B+hR_0)^2\\
 &\quad+\bigl[R_0-2c(\mathcal B+hR_0)(q-hA_0)\bigr]w\\
 &\quad-\bigl[A_0+c(q-hA_0)^2\bigr]w^2.
 \end{aligned}
\]
Its maximum over the whole real line is negative precisely under
\eqref{eq:enhanced-master}; the cancellation is
\[
 A_0(\mathcal B+hR_0)+R_0(q-hA_0)=A_0\mathcal B+qR_0.
\]
\end{proof}

This is the central dimension-six refinement.  The dimension-seven reduction
kept only the square-root term in \eqref{eq:penalized-energy}.  In dimension
six, the terms $hv$ and $qw$ supply exactly the missing coercivity.

\section{Exact two-sided Hall certificates}

Up to transposition, the six pairs are
\[
 (a,b)=(1,1),(1,2),(1,3),(1,4),(2,2),(2,3).
 \tag{5.1}\label{eq:six-pairs}
\]
Both excesses cannot be nonpositive, because then
$p\tau=e+f-x\le0$.  It remains to cover the three sign quadrants
\[
 (++),\qquad(-+),\qquad(+-)
\]
on the rational square $[0,T_0]^2$ after replacing a negative excess by the
negative of its square variable.

All quantities in the alternatives below are rational functions.  Put
\[
 Q_0=(m-e)(\ell-f)\omega,
 \qquad
 \widetilde{\mathcal B}=Q_0\mathcal B,
 \tag{5.2}\label{eq:Btilde}
\]
and
\[
 \widetilde h=Q_0h
 =(T+d)(m-e)(\ell-f)
 -p\omega\bigl(e_+(\ell-f)+f_+(m-e)\bigr).
 \tag{5.3}\label{eq:htilde}
\]
The positive-denominator cleared master polynomial is
\[
 \begin{aligned}
 \mathcal M_{\lambda,\theta}
 &=4\lambda\theta
 \bigl(\widetilde{\mathcal B}+\widetilde hR_0\bigr)
 \bigl(A_0\widetilde{\mathcal B}+qR_0Q_0\bigr)\\
 &\qquad-R_0^2Q_0^2.
 \end{aligned}
 \tag{5.4}\label{eq:Mpoly}
\]

Exact Bernstein subdivision proves that every relevant point satisfies one of
\[
 K<0,
 \qquad A_0-K>0,
 \qquad
 \widetilde{\mathcal B}>0\ \text{ and }\
 \mathcal M_{\lambda,\theta}>0.
 \tag{5.5}\label{eq:alternatives}
\]
In a mixed-sign quadrant, boxes on which $e+f\le0$ are discarded as
noncritical.  The first alternative in \eqref{eq:alternatives} contradicts
$H=K-u\ge0$; the second contradicts the permanent budget; and the third is
Lemma \ref{lem:penalized-triangle}.

For all pairs except $(2,3)$, the global coefficient $\theta=7/20$ is enough.
For $(2,3)$, a box uses the largest energy-ladder coefficient whose threshold
$A_*$ is certified by a Bernstein lower bound for $A_0$.

The exact subdivision audit is as follows.  The last three columns record
direct, stationarity, and noncritical leaves.

\begin{center}
\small
\begin{tabular}{c|c|r|r|r|r|r}
\toprule
$(a,b)$&quadrant&nodes&depth&direct&stationarity&noncritical\\
\midrule
$(1,1)$&$++$&21&10&5&6&0\\
       &$-+$&13&6&7&0&0\\
       &$+-$&11&5&4&0&2\\
\midrule
$(1,2)$&$++$&21&9&5&6&0\\
       &$-+$&9&4&5&0&0\\
       &$+-$&3&1&2&0&0\\
\midrule
$(1,3)$&$++$&39&8&11&9&0\\
       &$-+$&11&5&5&1&0\\
       &$+-$&3&1&2&0&0\\
\midrule
$(1,4)$&$++$&57&11&13&16&0\\
       &$-+$&29&8&9&5&1\\
       &$+-$&65&14&18&14&1\\
\midrule
$(2,2)$&$++$&23&8&6&6&0\\
       &$-+$&5&2&3&0&0\\
       &$+-$&3&1&2&0&0\\
\midrule
$(2,3)$&$++$&131&15&16&50&0\\
       &$-+$&23&8&7&5&0\\
       &$+-$&47&10&12&11&1\\
\bottomrule
\end{tabular}
\end{center}

Every subdivision is the exact rational de Casteljau transformation.  Thus the
audit proves positivity on continua, not on a numerical mesh.  This excludes
all two-sided critical Hall pairs.

\section{One-sided Hall cuts}

It remains to exclude a one-sided critical pair.  The stationarity identity,
column-entropy permanent bound, and entropy dichotomy from the dimensions
$7$--$10$ argument apply without change.  Only the exact constants improve.

The first row of the energy ladder gives, for
$y_i=1-1/r_i$,
\[
 \rho>\frac7{20}\sum_i y_i^2.
 \tag{6.1}\label{eq:row-energy6}
\]
For a one-sided pair, the subset estimate gives
\[
 \tau^2\le\frac{3\gamma}{1-\gamma}
 =\frac{15}{319}<T_0^2.
 \tag{6.2}\label{eq:one-sided-tau6}
\]
Let
\[
 \kappa_6=v_5G(5)=\frac{6144}{390625},
 \qquad
 s_0=2\log\frac{\kappa_6}{\gamma}.
\]

In the small-entropy regime, the exact sufficient certificate is
\[
 \frac7{20}(1-\gamma)^2\gamma-\frac{36}{2}\kappa_6^2
 =\frac{16250344917716167}
 {20759414062500000000}>0.
 \tag{6.3}\label{eq:small6}
\]

In the large-entropy regime, the exact inequality before squaring is
\[
 (1-\gamma)\sqrt{\frac{7\gamma}{10}}
 (1-\tau)^2
 >
 \gamma\frac{1-(1-\tau)^6}{\tau}.
 \tag{6.4}\label{eq:large-before-square}
\]
It is enough to prove, for $0\le t\le217/1000$,
\[
 \begin{aligned}
 &\frac7{10}(1-\gamma)^2(1-t)^4\\
 &\quad-\gamma
 (6-15t+20t^2-15t^3+6t^4-t^5)^2>0.
 \end{aligned}
 \tag{6.5}\label{eq:large6-poly}
\]
Every Bernstein coefficient of \eqref{eq:large6-poly} on this interval is
positive; the minimum is
\[
 \frac{3951217684118237014875556869551}
 {64800000000000000000000000000000}>0.
 \tag{6.6}\label{eq:large6-min}
\]
The existing small/large entropy argument therefore excludes every one-sided
critical pair.

\section{Completion}

If $\tau=0$, Hall domination makes a boundary maximizer itself doubly
stochastic.  The one-zero permanent bound gives
\[
 \per(A)\ge\kappa_6>\gamma_6,
 \qquad
 \kappa_6-\gamma_6=\frac{37531}{126562500}>0,
\]
contradicting the deficit budget.

If $\tau>0$, Section 5 excludes all two-sided critical pairs and Section 6
excludes all one-sided pairs.  Hence a global maximizer cannot lie on the
boundary.  The arbitrary-dimensional positive-maximizer theorem then gives
$A=J_6/6$, proving Theorem \ref{thm:n6}.

\begin{remark}
The essential new point is not a larger scalar constant.  The critical core
has two independent resources: marginal entropy raises $H_0$, while departure
from the groupwise AM--GM envelopes costs row/column product.  In the earlier
reduction both resources were collapsed into the single estimate
$\mathcal B\le\sqrt{su/\theta}$.  Keeping the terms $qw$ and $hv$ converts the
feasible region into a two-variable triangle whose exact quadratic maximum can
be eliminated.  That is the additional precision required at $n=6$.
\end{remark}

\end{document}
